\documentclass[10pt,a4paper]{amsart}
\usepackage[margin=19mm]{geometry}
\usepackage{amsmath,amssymb,mathtools}
\usepackage[scr=rsfs,cal=euler]{mathalfa}
\usepackage{xfrac}
\usepackage[unicode,hidelinks,bookmarks=false]{hyperref}
\newcommand{\ie}{i.e.\ }
\newcommand{\cf}{cf.\ }
\newcommand{\resp}{respectively\ }
\newcommand{\Subsection}{\subsection}
\providecommand{\nobreakdash}{\nobreak\hbox{-}}
\setlength{\emergencystretch}{2em}
\multlinegap0pt
\newtheorem{Th}{Theorem}
\newtheorem{Lemma}[Th]{Lemma}
\newtheorem{Cor}[Th]{Corollary}
\title[Theorem 1.4]{The largest automorphism group of del Pezzo surface of degree $2$ without points (Theorem 1.4)}
\author{Anastasia V. Vikulova}
\date{Source: arXiv:2501.16267v2 (CC BY 4.0)}
\begin{document}
\maketitle

\noindent\emph{Source and licence.} The largest automorphism group of del Pezzo surface of degree $2$ without points, by Anastasia V. Vikulova. Version arXiv:2501.16267v2
(CC BY 4.0), distributed by arXiv under the Creative Commons Attribution 4.0 International licence,
http://creativecommons.org/licenses/by/4.0/, as recorded in the arXiv licence metadata for this exact
version and shown on the abstract page https://arxiv.org/abs/2501.16267v2. Full licence notice:
\texttt{licenses/arxiv-2501.16267-CC-BY-4.0.txt}.\par\smallskip

\noindent\textit{Editorial scope.} Counted unit: the article's Theorem 1.4 (source0.tex lines 129--159, printed as Theorem 1.4): the surface $S$ of degree $2$ over $\mathbb{Q}(\sqrt{-7})$ given by the displayed quartic-plus-square equation in $\mathbb{P}(2,1,1,1)$ is smooth, has $\mathrm{Aut}(S)\simeq \mathrm{PSL}_2(\mathbb{F}_7)\times\mathbb{Z}/2\mathbb{Z}$ and $\mathrm{Pic}(S)\simeq\mathbb{Z}$, and has no rational points. The article's complete original proof of it is delivered with it (source0.tex lines 415--500, one \texttt{proof} environment of 86 lines carrying the article's own optional title, adjacent to the statement in this document, with no gap and no deferral). No mathematics is written, repaired or re-derived: the statement and the proof are the article's own lines, in the article's own notation, and no retained line is substituted, re-flowed or omitted, so nothing is removed and the package carries no omission record. Retained uncounted as the source-local context the proof needs, each exactly the statement of the result it invokes and nothing else: the completion lemma over $\mathbb{Q}(\sqrt{-7})$ (lines 205--214, printed as Lemma 2.3), the corollary reducing the surface equation modulo $64$ (lines 326--332, printed as Corollary 2.6, source label \texttt{cor}) and the centralizer lemma inside $W(E_7)$ (lines 355--357, printed as Lemma 2.7, source label \texttt{lemma:centralizer}). Every cross-reference of every retained line resolves inside this document.

Bibliography: the retained statement and proof contain exactly two cited works, and their entries are transcribed verbatim from the article's own bibliography block in source0.tex (the entries on lines 584 and 588--590) and delivered as a two-entry bibliography with short editorial labels; the mapping is recorded in the item's reference mapping. No other bibliography entry is transcribed, nothing is quoted from any work that the retained lines do not cite, and no retained line contains any other citation.

Editorial formatting only, in addition: an AMS \texttt{amsart} preamble that restates the article's own declarations used by the retained lines, namely the environments \texttt{Th}, \texttt{Lemma} and \texttt{Cor} on one shared counter exactly as in the article but without the article's subsection-based prefix. The retained environments therefore print in this document as Lemma 1, Corollary 2, Lemma 3 and Theorem 4 in order of appearance, whereas the article prints the same four items as Lemma 2.3, Corollary 2.6, Lemma 2.7 and Theorem 1.4; the displayed equations of the retained lines print as (1) to (5) in order of appearance, whereas the article numbers its equations per section.

The complete native source of the article as distributed in the arXiv e-print for this exact version is bundled unchanged as \texttt{sources/172262/source0.tex}: it is byte identical to the e-print file \texttt{Brauer-ManindP2.tex} except that the pipeline's mechanical concatenation prepends one header comment line naming it. The article uses the AMS \texttt{amsart} class; no publisher class or style file is shipped in the e-print and none is redistributed or used.
\begin{Lemma}\label{lemma:completeQ2}
There is a valuation $|\cdot|_v$ on  $\mathbb{Q}(\sqrt{-7})$ such that the completion~$\mathbb{Q}(\sqrt{-7})_v$ of   $\mathbb{Q}(\sqrt{-7})$ with respect to this valuation is isomorphic to~$\mathbb{Q}_2.$ Moreover, one can choose an isomorphism $\theta:\mathbb{Q}(\sqrt{-7})_{v} \xrightarrow{\sim} \mathbb{Q}_2$ so that
$$
\theta(\sqrt{-7}) \equiv 181 \mod 128
$$
\noindent and 
$$
\theta\left(\frac{3}{2}(1-\sqrt{-7})\right) \equiv -14 \mod 64.
$$ 
\end{Lemma}

\begin{Cor}\label{cor}
The equation 
\begin{equation}\label{eq:equationincorolary}
w^2+(x^4+y^4+z^4)+14(x^2y^2+x^2z^2+y^2z^2)=0
\end{equation}
\noindent does not have  non-zero solutions in $\mathbb{Q}_2.$
\end{Cor}

\begin{Lemma}\label{lemma:centralizer}
The group $\mathrm{PSL}_2(\mathbb{F}_7)$ lies in $W(\mathrm{E}_7),$ it is a simple group and its centralizer in  $W(\mathrm{E}_7)$ is isomorphic to $\mathbb{Z}/2\mathbb{Z}.$
\end{Lemma}

\begin{Th}\label{th:nopoints}
Let $S$ be a surface defined by the equation
$$
w^2+(x^4+y^4+z^4)-\frac{3}{2}(1-\sqrt{-7})(x^2y^2+x^2z^2+y^2z^2)=0
$$
\noindent over $\mathbb{Q}(\sqrt{-7})$ in the weighted projective space $\mathbb{P}(2,1,1,1)$ with  homogeneous coordinates~$w,x,y,z$ such that $\deg(w)=2$ and $\deg(x)=\deg(y)=\deg(z)=1.$ Then the following holds:
\begin{enumerate}
\renewcommand\labelenumi{\rm (\roman{enumi})}
\renewcommand\theenumi{\rm (\roman{enumi})}


\item\label{conditiondp2}
$S$ is a smooth del Pezzo surface of degree $2;$


\item\label{conditionAut}
one has $\mathrm{Aut}(S)\simeq \mathrm{PSL}_2(\mathbb{F}_7) \times \mathbb{Z}/2\mathbb{Z};$

\item\label{conditionPic}
$\mathrm{Pic}(S) \simeq \mathbb{Z};$

\item\label{conditionnopoint}
$S$ does not have rational points.

% over $\mathbb{Q}(\sqrt{-7}).$


\end{enumerate}


\end{Th}

\begin{proof}[Proof of Theorem \ref{th:nopoints}]
Let us prove \ref{conditiondp2}. The smoothness of the surface $S$ can be proved by direct computation of partial derivatives. By~\cite[Proposition 6.3.6]{Dolgachev} the surface $S$ is a del Pezzo surface of degree $2.$



Let us prove \ref{conditionAut}. It is well-known that the anticanonical linear system of $S$ gives a double cover $\phi_{|-K_S|}:S \to \mathbb{P}^2$ which is ramified over a smooth quartic curve $C$. In our case the equation of the quartic $C$ is 
\begin{equation}\label{eq:quarticcurve}
(x^4+y^4+z^4)-\frac{3}{2}(1-\sqrt{-7})(x^2y^2+x^2z^2+y^2z^2)=0
\end{equation}
\noindent in $\mathbb{P}^2$ with homogeneous coordinates $x,y,z.$  The double cover $\phi_{|-K_S|}$ gives us the following exact sequence
\begin{equation}\label{eq:exactsequensecanonical}
0 \to \mathbb{Z}/2\mathbb{Z} \to \mathrm{Aut}(S) \to \mathrm{Aut}(\mathbb{P}^2),
\end{equation}
\noindent where $\mathbb{Z}/2\mathbb{Z}$ is generated by the Geiser involution~$\gamma$ which is a deck transformation of the cover  $\phi_{|-K_S|}$ and  acts on $S$ as 
$$
\gamma:(w,x,y,z) \mapsto (-w,x,y,z).
$$
\noindent From \eqref{eq:exactsequensecanonical} we obtain
\begin{equation}\label{eq:aut}
\mathrm{Aut}(S)/\left(\mathbb{Z}/2\mathbb{Z}\right) \simeq \mathrm{Aut}(C) \simeq \mathrm{PSL}_2(\mathbb{F}_7),
\end{equation}
\noindent where the last isomorphism follows  from~\cite[Chapter 1]{Elkies}. Note that we have the natural embedding $\mathrm{PSL}_2(\mathbb{F}_7) \subset \mathrm{Aut}(S).$  Since  the Geiser involution~\mbox{$\gamma \in \mathrm{Aut}(S)$} and $\mathrm{PSL}_2(\mathbb{F}_7) \subset \mathrm{Aut}(S)$ commute with each other, we get the desired isomorphism
$$
\mathrm{PSL}_2(\mathbb{F}_7) \times \mathbb{Z}/2\mathbb{Z} \simeq \mathrm{Aut}(S).
$$














Let us prove \ref{conditionPic}. To prove that $\mathrm{Pic}(S) \simeq \mathbb{Z}$ one should prove that  
\begin{equation}\label{eq:picgaloisinvariant}
\mathrm{Pic}(\bar{S})^{\mathrm{Gal}(\bar{\mathbb{Q}}/\mathbb{Q}(\sqrt{-7}))} \simeq \mathbb{Z},
\end{equation}
\noindent where $\bar{S}=S \times_{\mathrm{Spec}(\mathbb{Q}(\sqrt{-7}))}\mathrm{Spec}(\bar{\mathbb{Q}}),$ because we have the obvious inclusion 
$$
\mathrm{Pic}(S) \subseteq \mathrm{Pic}(\bar{S})^{\mathrm{Gal}(\bar{\mathbb{Q}}/\mathbb{Q}(\sqrt{-7}))}.
$$
\noindent To prove \eqref{eq:picgaloisinvariant} note that the image of the Galois group $\mathrm{Gal}(\bar{\mathbb{Q}}/\mathbb{Q}(\sqrt{-7}))$ in the Weyl group $W(\mathrm{E}_7)$ and the automorphism group $\mathrm{Aut}(S),$ which lies in $W(\mathrm{E}_7)$ by~\cite[Corollary 8.2.40]{Dolgachev}, commute with each other. So we need to consider the centralizer of~\mbox{$\mathrm{Aut}(S) \simeq \mathrm{PSL}_2(\mathbb{F}_7) \times \mathbb{Z}/2\mathbb{Z}$} in $W(\mathrm{E}_7).$ By Lemma \ref{lemma:centralizer} it is isomorphic to~$\mathbb{Z}/2\mathbb{Z}$ and it is obvious that it is generated by the Geiser involution. This means that the image of the Galois group $\mathrm{Gal}(\bar{\mathbb{Q}}/\mathbb{Q}(\sqrt{-7}))$ in the Weyl group $W(\mathrm{E}_7)$ is either trivial, or isomorphic to $\mathbb{Z}/2\mathbb{Z} \simeq \langle \gamma \rangle.$ 

To prove that it is not trivial it is enough to show that that there is an exceptional curve on $S$ which is not defined over $\mathbb{Q}(\sqrt{-7}).$ According to the basic theory of smooth del Pezzo surfaces of degree $2$ (see, for instance,~\cite[\S 8.7.1]{Dolgachev}) the surface $S$ has $56$ lines which are the inverse image of $28$ bitangents of the branch curve $C.$ It is not hard to see that the line in $\mathbb{P}^2$ defined by the equation
$$
x+y+z=0
$$
is a bitangent of $C.$ So the inverse image of this bitangent are two lines which are defined by the equations
$$
w \pm i\left(\frac{3+\sqrt{-7}}{2} \right)^2(x^2+xy+y^2)=0.
$$
\noindent They are not defined over $\mathbb{Q}(\sqrt{-7}),$ hence, the image of  $\mathrm{Gal}(\bar{\mathbb{Q}}/\mathbb{Q}(\sqrt{-7}))$ in the Weyl group $W(\mathrm{E}_7)$ is isomorphic to $\mathbb{Z}/2\mathbb{Z} \simeq \langle \gamma \rangle.$ As it is generated by the Geiser involution which  interchanges the two lines in the inverse image of
any bitangent of $C$, we get
$$ 
\mathrm{Pic}(\bar{S})^{\mathrm{Gal}(\bar{\mathbb{Q}}/\mathbb{Q}(\sqrt{-7}))} \simeq \mathbb{Z}.
$$







Let us prove \ref{conditionnopoint}. To prove that $S$ does not have a rational point over the field~$\mathbb{Q}(\sqrt{-7})$ it is enough to prove that  it does not have a rational point over some completion of $\mathbb{Q}(\sqrt{-7}).$ Let us consider the valuation $|\cdot|_2$ on $\mathbb{Q}.$  Let $|\cdot|_v$ be one of the extensions of $|\cdot|_2$ on $\mathbb{Q}(\sqrt{-7})$ (which was defined in Lemma~\ref{lemma:completeQ2}). Then by Lemma~\ref{lemma:completeQ2} for the completion~$\mathbb{Q}(\sqrt{-7})_v$ of $\mathbb{Q}(\sqrt{-7})$ with respect to $|\cdot|_v$ we have an isomorphism
$$
\theta:\mathbb{Q}(\sqrt{-7})_v  \xrightarrow{\sim} \mathbb{Q}_2.
$$ 
This isomorphism satisfies
$$
\theta\left(\frac{3}{2}(1-\sqrt{-7})\right) \equiv -14 \mod 64.
$$ 

\noindent So to prove that  $S$ does not have a rational point over the field $\mathbb{Q}_2$ it is enough to prove that
$$
w^2+(x^4+y^4+z^4)+14(x^2y^2+x^2z^2+y^2z^2)=0
$$
\noindent does not have a solution over $\mathbb{Q}_2.$   However, by Corollary \ref{cor} it does not have a non-zero solution. Quod erat demonstrandum.


\end{proof}

\begin{thebibliography}{99}
\bibitem[Dol12]{Dolgachev} I.~V.~Dolgachev, \emph{Classical algebraic geometry. A modern view}, Cambridge University Press, Cambridge (2012).
\bibitem[Elk01]{Elkies} N.~D.~Elkies, \emph{The Klein quartic in number theory}, In S.~Levi (ed.), \emph{The eightfold way: the beauty of Klein's quartic curve},   Mathematical Sciences Research Institute publications, Cambridge University Press, Cambridge, \textbf{35} (2001), 51--102.
\end{thebibliography}
\end{document}
